\section{左右の値を組にする}\label{節:奇偶の対称性}

\begin{definition}{偶関数と奇関数}
原点に関して対称な範囲で定義された関数 $f$ について，すべての $x$ で
\[
f(-x)=f(x)
\]
が成り立つとき，$f$ を偶関数という。また，すべての $x$ で
\[
f(-x)=-f(x)
\]
が成り立つとき，$f$ を奇関数という。
\end{definition}

偶関数のグラフは $y$ 軸について対称であり，奇関数のグラフは原点について対称である。$x^2$ と $\cos x$ は偶関数，$x^3$ と $\sin x$ は奇関数である。ただし，式の中に $\sin$ があれば奇関数になるわけではない。$\sin(x^2)$ では $x$ を $-x$ に変えても $x^2$ が変わらないので偶関数である。式全体に $-x$ を代入することが，確実な判定になる。

$a>0$ とし，$[-a,a]$ の積分を零で分ける。左半分で $x=-t$ と置くと，$x=-a,0$ に対して $t=a,0$ であり，$\dd x=-\dd t$ だから，
\[
\int_{-a}^{0}f(x)\dd x
 =\int_a^0 f(-t)(-\dd t)
 =\int_0^a f(-t)\dd t.
\]
積分変数の文字をそろえれば，
\begin{equation}\label{式:左右の積分}
\int_{-a}^a f(x)\dd x
 =\int_0^a\{f(x)+f(-x)\}\dd x
\end{equation}
となる。これは偶関数でも奇関数でもない場合にも使える式である。偶関数なら中括弧が $2f(x)$ になり，奇関数なら零になる。

\begin{theorem}{対称区間での偶関数・奇関数の積分}
$a>0$ とする。$f$ が偶関数なら
\[
\int_{-a}^a f(x)\dd x=2\int_0^a f(x)\dd x,
\]
$f$ が奇関数なら
\[
\int_{-a}^a f(x)\dd x=0
\]
である。
\end{theorem}

図\ref{図:奇偶の積分}の右図では，左右の部分の幾何学的な面積は等しい。しかし左側の定積分は負，右側は正である。奇関数の積分が零という結論は，図形の面積が零という意味ではない。面積を求めるなら絶対値を積分する必要がある。

\begin{figure}[htbp]
\centering
\begin{tikzpicture}[scale=1.05,>=Stealth]
\begin{scope}[xshift=-3.1cm]
\fill[blue!12] plot[domain=-1.35:1.35,samples=50](\x,{0.65*\x*\x})--(1.35,0)--(-1.35,0)--cycle;
\draw[->](-1.7,0)--(1.8,0)node[right]{$x$};
\draw[->](0,-0.2)--(0,1.7)node[above]{$y$};
\draw[blue!70!black,thick] plot[domain=-1.45:1.45,samples=60](\x,{0.65*\x*\x});
\draw[dashed](-1.35,0)--(-1.35,1.185);
\draw[dashed](1.35,0)--(1.35,1.185);
\node[below]at(-1.35,0){$-a$};\node[below]at(1.35,0){$a$};
\node at(0,-0.65){偶関数：左右を加える};
\end{scope}
\begin{scope}[xshift=1.8cm]
\fill[blue!12](0,0)--plot[domain=0:1.35,samples=40](\x,{0.55*\x*\x*\x})--(1.35,0)--cycle;
\fill[orange!18](0,0)--plot[domain=0:-1.35,samples=40](\x,{0.55*\x*\x*\x})--(-1.35,0)--cycle;
\draw[->](-1.7,0)--(1.8,0)node[right]{$x$};
\draw[->](0,-1.55)--(0,1.7)node[above]{$y$};
\draw[blue!70!black,thick] plot[domain=-1.4:1.4,samples=60](\x,{0.55*\x*\x*\x});
\draw[dashed](-1.35,0)--(-1.35,-1.353);
\draw[dashed](1.35,0)--(1.35,1.353);
\node[above]at(-1.35,0){$-a$};\node[below]at(1.35,0){$a$};
\node at(0,-1.95){奇関数：符号付きの量が打ち消す};
\end{scope}
\end{tikzpicture}
\caption{対称性と，定積分の符号の関係。曲線はそれぞれ $x^2$ 型，$x^3$ 型である。}\label{図:奇偶の積分}
\end{figure}

\begin{example}{混じった式から奇関数の部分を除く}
次の定積分を求めよ。
\[
I=\int_{-2}^{2}\left(x^3\e^{-x^2}+2x^2+1\right)\dd x.
\]
\end{example}

$x^3\e^{-x^2}$ の原始関数を先に探すこともできるが，区間が $[-2,2]$ なので，まず $x$ を $-x$ に変えてみる。$x^3$ の符号だけが変わり，$\e^{-x^2}$ は変わらない。したがって積 $x^3\e^{-x^2}$ は奇関数で，その積分は零である。残った $2x^2+1$ は偶関数だから，
\[
I=2\int_0^2(2x^2+1)\dd x
 =2\left[\frac{2x^3}{3}+x\right]_0^2
 =\frac{44}{3}.
\]
ここでは式全体が偶関数か奇関数かを決める必要はなかった。和を分けて，対称性が分かる部分にだけ使えばよい。

さらに，一般の $f$ に対して
\[
f_{\mathrm e}(x)=\frac{f(x)+f(-x)}{2},\qquad
f_{\mathrm o}(x)=\frac{f(x)-f(-x)}{2}
\]
と置くと，$f=f_{\mathrm e}+f_{\mathrm o}$ であり，前者は偶関数，後者は奇関数になる。よって対称区間の定積分に寄与するのは $f_{\mathrm e}$ の部分だけである。これは新しい積分公式を覚えるための分解ではなく，式\eqref{式:左右の積分}が何を選び出しているかを示す分解である。

